% Môn: Vật lý
% Lớp: 12
% Chương: Chương II. Khí lí tưởng
% Bài: L12C2 Bài 13. Bài tập về khí lí tưởng · Nhận diện quá trình nhiệt động từ đồ thị và dữ kiện
% ID: L12C1B2-01-DS
% Mức: NB
% ID: L12C1B2-01-DS


% ID: L12C1B2-11-DS
% Mức: NB
% ID: L12C1B2-11-DS
% ID: L12C1B2-11-DS
% Mức: NB
% ID: L12C1B2-11-DS
% ID: L12C1B2-11-DS
% Mức: NB
% ID: L12C2B13-61-DS
%=== Câu 28 ===%
\dangbt{Nhận diện quá trình nhiệt động từ đồ thị và dữ kiện}
\begin{ex}
% ID: L12C2B13-61-DS
% Mức: NB
Xét các phát biểu sau về dạng “Nhận diện quá trình nhiệt động từ đồ thị và dữ kiện”
\choiceTF
{\True Khí nhận nhiệt và thực hiện công có ΔU=Q-A_do_khí_thực_hiện.}
{Khí nhận nhiệt thì nội năng luôn tăng đúng bằng nhiệt lượng nhận được.}
{\True Khi giải dạng “Nhận diện quá trình nhiệt động từ đồ thị và dữ kiện”, cần xác định đúng trạng thái đầu, trạng thái cuối và quy ước dấu hoặc điều kiện áp dụng.}
{Có thể bỏ qua mọi dữ kiện về khối lượng, nhiệt độ và hiệu suất trong mọi bài toán thuộc dạng này.}
\loigiai{
\begin{itemchoice}
\item (Đúng) Khí nhận nhiệt và thực hiện công có ΔU=Q-A_do_khí_thực_hiện. (Vì): _do_khí_thực_hiện. b) S: Khí nhận nhiệt thì nội năng luôn tăng đúng bằng nhiệt lượng nhận được. c) Đ: Khi giải dạng “Nhận diện quá trình nhiệt động từ đồ thị và dữ kiện”, cần xác định đúng trạng thái đầu, trạng thái cuối và quy ước dấu hoặc điều kiện áp dụng. d) S: Có thể bỏ qua mọi dữ kiện về khối lượng, nhiệt độ và hiệu suất trong mọi bài toán thuộc dạng này. Kết luận dựa trên định nghĩa, điều kiện áp dụng và quy ước trong SGK KNTT
\item (Sai) Khí nhận nhiệt thì nội năng luôn tăng đúng bằng nhiệt lượng nhận được.
\item (Đúng) Khi giải dạng “Nhận diện quá trình nhiệt động từ đồ thị và dữ kiện”, cần xác định đúng trạng thái đầu, trạng thái cuối và quy ước dấu hoặc điều kiện áp dụng.
\item (Sai) Có thể bỏ qua mọi dữ kiện về khối lượng, nhiệt độ và hiệu suất trong mọi bài toán thuộc dạng này.
\end{itemchoice}
}
\end{ex}

% ID: L12C1B2-12-DS
% Mức: TH
% ID: L12C1B2-12-DS
% ID: L12C1B2-12-DS
% Mức: TH
% ID: L12C1B2-12-DS
% ID: L12C1B2-12-DS
% Mức: TH
% ID: L12C2B13-62-DS
%=== Câu 30 ===%
\begin{ex}
% ID: L12C2B13-62-DS
% Mức: TH
Xét các phát biểu sau về dạng “Nhận diện quá trình nhiệt động từ đồ thị và dữ kiện”
\choiceTF
{Khí nhận nhiệt thì nội năng luôn tăng đúng bằng nhiệt lượng nhận được.}
{\True Khi giải dạng “Nhận diện quá trình nhiệt động từ đồ thị và dữ kiện”, cần xác định đúng trạng thái đầu, trạng thái cuối và quy ước dấu hoặc điều kiện áp dụng.}
{Có thể bỏ qua mọi dữ kiện về khối lượng, nhiệt độ và hiệu suất trong mọi bài toán thuộc dạng này.}
{\True Khí nhận nhiệt và thực hiện công có ΔU=Q-A_do_khí_thực_hiện.}
\loigiai{
\begin{itemchoice}
\item (Sai) Khí nhận nhiệt thì nội năng luôn tăng đúng bằng nhiệt lượng nhận được. (Vì): _do_khí_thực_hiện. Kết luận dựa trên định nghĩa, điều kiện áp dụng và quy ước trong SGK KNTT
\item (Đúng) Khi giải dạng “Nhận diện quá trình nhiệt động từ đồ thị và dữ kiện”, cần xác định đúng trạng thái đầu, trạng thái cuối và quy ước dấu hoặc điều kiện áp dụng.
\item (Sai) Có thể bỏ qua mọi dữ kiện về khối lượng, nhiệt độ và hiệu suất trong mọi bài toán thuộc dạng này.
\item (Đúng) Khí nhận nhiệt và thực hiện công có ΔU=Q-A_do_khí_thực_hiện.
\end{itemchoice}
}
\end{ex}

% ID: L12C1B2-13-DS
% Mức: TH
% ID: L12C1B2-13-DS
% ID: L12C1B2-13-DS
% Mức: TH
% ID: L12C1B2-13-DS
% ID: L12C1B2-13-DS
% Mức: TH
% ID: L12C2B13-63-DS
%=== Câu 32 ===%
\begin{ex}
% ID: L12C2B13-63-DS
% Mức: TH
Xét các phát biểu sau về dạng “Nhận diện quá trình nhiệt động từ đồ thị và dữ kiện”
\choiceTF
{Khí nhận nhiệt thì nội năng luôn tăng đúng bằng nhiệt lượng nhận được.}
{\True Khi giải dạng “Nhận diện quá trình nhiệt động từ đồ thị và dữ kiện”, cần xác định đúng trạng thái đầu, trạng thái cuối và quy ước dấu hoặc điều kiện áp dụng.}
{Có thể bỏ qua mọi dữ kiện về khối lượng, nhiệt độ và hiệu suất trong mọi bài toán thuộc dạng này.}
{\True Khí nhận nhiệt và thực hiện công có ΔU=Q-A_do_khí_thực_hiện.}
\loigiai{
\begin{itemchoice}
\item (Sai) Khí nhận nhiệt thì nội năng luôn tăng đúng bằng nhiệt lượng nhận được. (Vì): _do_khí_thực_hiện. Kết luận dựa trên định nghĩa, điều kiện áp dụng và quy ước trong SGK KNTT
\item (Đúng) Khi giải dạng “Nhận diện quá trình nhiệt động từ đồ thị và dữ kiện”, cần xác định đúng trạng thái đầu, trạng thái cuối và quy ước dấu hoặc điều kiện áp dụng.
\item (Sai) Có thể bỏ qua mọi dữ kiện về khối lượng, nhiệt độ và hiệu suất trong mọi bài toán thuộc dạng này.
\item (Đúng) Khí nhận nhiệt và thực hiện công có ΔU=Q-A_do_khí_thực_hiện.
\end{itemchoice}
}
\end{ex}

% ID: L12C1B2-13-TL
% Mức: NB
% ID: L12C1B2-13-TL
% ID: L12C1B2-13-TL
% Mức: NB
% ID: L12C1B2-13-TL
% ID: L12C1B2-13-TL
% Mức: NB
% ID: L12C2B13-64-TL
%=== Câu 33 ===%
\begin{bt}
% ID: L12C2B13-64-TL
% Mức: NB
Trình bày cách nhận biết và phương pháp giải dạng “Nhận diện quá trình nhiệt động từ đồ thị và dữ kiện”. Phân tích nhận định sau và sửa lại nếu sai: “Khí nhận nhiệt thì nội năng luôn tăng đúng bằng nhiệt lượng nhận được.”
\loigiai{
Trước hết xác định hiện tượng, trạng thái và các đại lượng đã cho. Nêu định nghĩa hoặc hệ thức phù hợp, đổi về đơn vị SI, áp dụng đúng điều kiện và kiểm tra ý nghĩa kết quả. Nhận định cần đối chiếu: Khí nhận nhiệt và thực hiện công có ΔU=Q-A_do_khí_thực_hiện. Vì vậy nhận định đã cho sai; phát biểu đúng là: Khí nhận nhiệt và thực hiện công có ΔU=Q-A_do_khí_thực_hiện.
}
\end{bt}

% ID: L12C1B2-14-DS
% Mức: VD
% ID: L12C1B2-14-DS
% ID: L12C1B2-14-DS
% Mức: VD
% ID: L12C1B2-14-DS
% ID: L12C1B2-14-DS
% Mức: VD
% ID: L12C2B13-65-DS
%=== Câu 35 ===%
\begin{ex}
% ID: L12C2B13-65-DS
% Mức: VD
Xét các phát biểu sau về dạng “Nhận diện quá trình nhiệt động từ đồ thị và dữ kiện” 
\choiceTF
{\True Khi giải dạng “Nhận diện quá trình nhiệt động từ đồ thị và dữ kiện”, cần xác định đúng trạng thái đầu, trạng thái cuối và quy ước dấu hoặc điều kiện áp dụng.}
{Có thể bỏ qua mọi dữ kiện về khối lượng, nhiệt độ và hiệu suất trong mọi bài toán thuộc dạng này.}
{\True Khí nhận nhiệt và thực hiện công có ΔU=Q-A_do_khí_thực_hiện.}
{Khí nhận nhiệt thì nội năng luôn tăng đúng bằng nhiệt lượng nhận được.}
\loigiai{
\begin{itemchoice}
\item (Đúng) Khi giải dạng “Nhận diện quá trình nhiệt động từ đồ thị và dữ kiện”, cần xác định đúng trạng thái đầu, trạng thái cuối và quy ước dấu hoặc điều kiện áp dụng. (Vì): _do_khí_thực_hiện. d) S: Khí nhận nhiệt thì nội năng luôn tăng đúng bằng nhiệt lượng nhận được. Kết luận dựa trên định nghĩa, điều kiện áp dụng và quy ước trong SGK KNTT
\item (Sai) Có thể bỏ qua mọi dữ kiện về khối lượng, nhiệt độ và hiệu suất trong mọi bài toán thuộc dạng này.
\item (Đúng) Khí nhận nhiệt và thực hiện công có ΔU=Q-A_do_khí_thực_hiện.
\item (Sai) Khí nhận nhiệt thì nội năng luôn tăng đúng bằng nhiệt lượng nhận được.
\end{itemchoice}
}
\end{ex}

% ID: L12C1B2-14-TLN
% Mức: NB
% ID: L12C1B2-14-TLN
% ID: L12C1B2-14-TLN
% Mức: NB
% ID: L12C1B2-14-TLN
% ID: L12C1B2-14-TLN
% Mức: NB
% ID: L12C2B13-66-TLN
%=== Câu 36 ===%
\begin{ex}
% ID: L12C2B13-66-TLN
% Mức: NB
Cho bốn nhận định: (1) Khí nhận nhiệt và thực hiện công có ΔU=Q-A_do_khí_thực_hiện. (2) Khí nhận nhiệt thì nội năng luôn tăng đúng bằng nhiệt lượng nhận được. (3) Cần kiểm tra đơn vị và điều kiện áp dụng khi xử lí dạng “Nhận diện quá trình nhiệt động từ đồ thị và dữ kiện”. (4) Có thể thay tùy ý nhiệt độ Celsius cho nhiệt độ tuyệt đối trong mọi công thức. Có bao nhiêu nhận định đúng? Trả lời bằng một số nguyên.
\shortans{2}
\loigiai{
Nhận định (1) và (3) đúng; (2) và (4) sai. Vậy có 2 nhận định đúng.
}
\end{ex}

% ID: L12C1B2-15-DS
% Mức: VDC
% ID: L12C1B2-15-DS
% ID: L12C1B2-15-DS
% Mức: VDC
% ID: L12C1B2-15-DS
% ID: L12C1B2-15-DS
% Mức: VDC
% ID: L12C2B13-67-DS
%=== Câu 38 ===%
\begin{ex}
% ID: L12C2B13-67-DS
% Mức: VDC
Xét các phát biểu sau về dạng “Nhận diện quá trình nhiệt động từ đồ thị và dữ kiện”
\choiceTF
{Có thể bỏ qua mọi dữ kiện về khối lượng, nhiệt độ và hiệu suất trong mọi bài toán thuộc dạng này.}
{\True Khí nhận nhiệt và thực hiện công có ΔU=Q-A_do_khí_thực_hiện.}
{Khí nhận nhiệt thì nội năng luôn tăng đúng bằng nhiệt lượng nhận được.}
{\True Khi giải dạng “Nhận diện quá trình nhiệt động từ đồ thị và dữ kiện”, cần xác định đúng trạng thái đầu, trạng thái cuối và quy ước dấu hoặc điều kiện áp dụng.}
\loigiai{
\begin{itemchoice}
\item (Sai) Có thể bỏ qua mọi dữ kiện về khối lượng, nhiệt độ và hiệu suất trong mọi bài toán thuộc dạng này. (Vì): _do_khí_thực_hiện. c) S: Khí nhận nhiệt thì nội năng luôn tăng đúng bằng nhiệt lượng nhận được. d) Đ: Khi giải dạng “Nhận diện quá trình nhiệt động từ đồ thị và dữ kiện”, cần xác định đúng trạng thái đầu, trạng thái cuối và quy ước dấu hoặc điều kiện áp dụng. Kết luận dựa trên định nghĩa, điều kiện áp dụng và quy ước trong SGK KNTT
\item (Đúng) Khí nhận nhiệt và thực hiện công có ΔU=Q-A_do_khí_thực_hiện.
\item (Sai) Khí nhận nhiệt thì nội năng luôn tăng đúng bằng nhiệt lượng nhận được.
\item (Đúng) Khi giải dạng “Nhận diện quá trình nhiệt động từ đồ thị và dữ kiện”, cần xác định đúng trạng thái đầu, trạng thái cuối và quy ước dấu hoặc điều kiện áp dụng.
\end{itemchoice}
}
\end{ex}

% ID: L12C1B2-15-TL
% Mức: TH
% ID: L12C1B2-15-TL
% ID: L12C1B2-15-TL
% Mức: TH
% ID: L12C1B2-15-TL
% ID: L12C1B2-15-TL
% Mức: TH
% ID: L12C2B13-68-TL
%=== Câu 39 ===%
\begin{bt}
% ID: L12C2B13-68-TL
% Mức: TH
Trình bày cách nhận biết và phương pháp giải dạng “Nhận diện quá trình nhiệt động từ đồ thị và dữ kiện”. Phân tích nhận định sau và sửa lại nếu sai: “Khí nhận nhiệt và thực hiện công có ΔU=Q-A_do_khí_thực_hiện.”
\loigiai{
Trước hết xác định hiện tượng, trạng thái và các đại lượng đã cho. Nêu định nghĩa hoặc hệ thức phù hợp, đổi về đơn vị SI, áp dụng đúng điều kiện và kiểm tra ý nghĩa kết quả. Nhận định cần đối chiếu: Khí nhận nhiệt và thực hiện công có ΔU=Q-A_do_khí_thực_hiện. Vì vậy nhận định đã cho đúng.
}
\end{bt}

% ID: L12C1B2-21-TN
% Mức: NB
% ID: L12C1B2-21-TN
% ID: L12C1B2-21-TN
% Mức: NB
% ID: L12C1B2-21-TN
% ID: L12C1B2-21-TN
% Mức: NB
% ID: L12C2B13-69-TN
%=== Câu 55 ===%
\begin{ex}
% ID: L12C2B13-69-TN
% Mức: NB
Câu 1 thuộc dạng “Nhận diện quá trình nhiệt động từ đồ thị và dữ kiện”. Phát biểu nào sau đây đúng?
\choice
{Khí nhận nhiệt thì nội năng luôn tăng đúng bằng nhiệt lượng nhận được.}
{Nhiệt lượng và nhiệt độ là hai đại lượng hoàn toàn giống nhau.}
{Mọi quá trình nhiệt học đều không phụ thuộc điều kiện thực hiện.}
{\True Khí nhận nhiệt và thực hiện công có ΔU=Q-A_do_khí_thực_hiện.}
\loigiai{
Phát biểu đúng là: Khí nhận nhiệt và thực hiện công có ΔU=Q-A_do_khí_thực_hiện. Các phương án còn lại trái với mô hình, định nghĩa hoặc hệ thức nhiệt học tương ứng.
}
\end{ex}

% ID: L12C1B2-22-TN
% Mức: NB
% ID: L12C1B2-22-TN
% ID: L12C1B2-22-TN
% Mức: NB
% ID: L12C1B2-22-TN
% ID: L12C1B2-22-TN
% Mức: NB
% ID: L12C2B13-70-TN
%=== Câu 57 ===%
\begin{ex}
% ID: L12C2B13-70-TN
% Mức: NB
Câu 5 thuộc dạng “Nhận diện quá trình nhiệt động từ đồ thị và dữ kiện”. Phát biểu nào sau đây đúng?
\choice
{Khí nhận nhiệt thì nội năng luôn tăng đúng bằng nhiệt lượng nhận được.}
{Nhiệt lượng và nhiệt độ là hai đại lượng hoàn toàn giống nhau.}
{Mọi quá trình nhiệt học đều không phụ thuộc điều kiện thực hiện.}
{\True Khí nhận nhiệt và thực hiện công có ΔU=Q-A_do_khí_thực_hiện.}
\loigiai{
Phát biểu đúng là: Khí nhận nhiệt và thực hiện công có ΔU=Q-A_do_khí_thực_hiện. Các phương án còn lại trái với mô hình, định nghĩa hoặc hệ thức nhiệt học tương ứng.
}
\end{ex}

% ID: L12C1B2-23-TN
% Mức: NB
% ID: L12C1B2-23-TN
% ID: L12C1B2-23-TN
% Mức: NB
% ID: L12C1B2-23-TN
% ID: L12C1B2-23-TN
% Mức: NB
% ID: L12C2B13-71-TN
%=== Câu 59 ===%
\begin{ex}
% ID: L12C2B13-71-TN
% Mức: NB
Câu 9 thuộc dạng “Nhận diện quá trình nhiệt động từ đồ thị và dữ kiện”. Phát biểu nào sau đây đúng?
\choice
{Khí nhận nhiệt thì nội năng luôn tăng đúng bằng nhiệt lượng nhận được.}
{Nhiệt lượng và nhiệt độ là hai đại lượng hoàn toàn giống nhau.}
{Mọi quá trình nhiệt học đều không phụ thuộc điều kiện thực hiện.}
{\True Khí nhận nhiệt và thực hiện công có ΔU=Q-A_do_khí_thực_hiện.}
\loigiai{
Phát biểu đúng là: Khí nhận nhiệt và thực hiện công có ΔU=Q-A_do_khí_thực_hiện. Các phương án còn lại trái với mô hình, định nghĩa hoặc hệ thức nhiệt học tương ứng.
}
\end{ex}

% ID: L12C1B2-24-TN
% Mức: TH
% ID: L12C1B2-24-TN
% ID: L12C1B2-24-TN
% Mức: TH
% ID: L12C1B2-24-TN
% ID: L12C1B2-24-TN
% Mức: TH
% ID: L12C2B13-72-TN
%=== Câu 61 ===%
\begin{ex}
% ID: L12C2B13-72-TN
% Mức: TH
Câu 2 thuộc dạng “Nhận diện quá trình nhiệt động từ đồ thị và dữ kiện”. Phát biểu nào sau đây đúng?
\choice
{Nhiệt lượng và nhiệt độ là hai đại lượng hoàn toàn giống nhau.}
{Mọi quá trình nhiệt học đều không phụ thuộc điều kiện thực hiện.}
{\True Khí nhận nhiệt và thực hiện công có ΔU=Q-A_do_khí_thực_hiện.}
{Khí nhận nhiệt thì nội năng luôn tăng đúng bằng nhiệt lượng nhận được.}
\loigiai{
Phát biểu đúng là: Khí nhận nhiệt và thực hiện công có ΔU=Q-A_do_khí_thực_hiện. Các phương án còn lại trái với mô hình, định nghĩa hoặc hệ thức nhiệt học tương ứng.
}
\end{ex}

% ID: L12C1B2-25-TN
% Mức: TH
% ID: L12C1B2-25-TN
% ID: L12C1B2-25-TN
% Mức: TH
% ID: L12C1B2-25-TN
% ID: L12C1B2-25-TN
% Mức: TH
% ID: L12C2B13-73-TN
%=== Câu 64 ===%
\begin{ex}
% ID: L12C2B13-73-TN
% Mức: TH
Câu 6 thuộc dạng “Nhận diện quá trình nhiệt động từ đồ thị và dữ kiện”. Phát biểu nào sau đây đúng?
\choice
{Nhiệt lượng và nhiệt độ là hai đại lượng hoàn toàn giống nhau.}
{Mọi quá trình nhiệt học đều không phụ thuộc điều kiện thực hiện.}
{\True Khí nhận nhiệt và thực hiện công có ΔU=Q-A_do_khí_thực_hiện.}
{Khí nhận nhiệt thì nội năng luôn tăng đúng bằng nhiệt lượng nhận được.}
\loigiai{
Phát biểu đúng là: Khí nhận nhiệt và thực hiện công có ΔU=Q-A_do_khí_thực_hiện. Các phương án còn lại trái với mô hình, định nghĩa hoặc hệ thức nhiệt học tương ứng.
}
\end{ex}

% ID: L12C1B2-26-TN
% Mức: TH
% ID: L12C1B2-26-TN
% ID: L12C1B2-26-TN
% Mức: TH
% ID: L12C1B2-26-TN
% ID: L12C1B2-26-TN
% Mức: TH
% ID: L12C2B13-74-TN
%=== Câu 67 ===%
\begin{ex}
% ID: L12C2B13-74-TN
% Mức: TH
Câu 10 thuộc dạng “Nhận diện quá trình nhiệt động từ đồ thị và dữ kiện”. Phát biểu nào sau đây đúng?
\choice
{Nhiệt lượng và nhiệt độ là hai đại lượng hoàn toàn giống nhau.}
{Mọi quá trình nhiệt học đều không phụ thuộc điều kiện thực hiện.}
{\True Khí nhận nhiệt và thực hiện công có ΔU=Q-A_do_khí_thực_hiện.}
{Khí nhận nhiệt thì nội năng luôn tăng đúng bằng nhiệt lượng nhận được.}
\loigiai{
Phát biểu đúng là: Khí nhận nhiệt và thực hiện công có ΔU=Q-A_do_khí_thực_hiện. Các phương án còn lại trái với mô hình, định nghĩa hoặc hệ thức nhiệt học tương ứng.
}
\end{ex}

% ID: L12C1B2-27-TN
% Mức: VD
% ID: L12C1B2-27-TN
% ID: L12C1B2-27-TN
% Mức: VD
% ID: L12C1B2-27-TN
% ID: L12C1B2-27-TN
% Mức: VD
% ID: L12C2B13-75-TN
%=== Câu 71 ===%
\begin{ex}
% ID: L12C2B13-75-TN
% Mức: VD
Câu 3 thuộc dạng “Nhận diện quá trình nhiệt động từ đồ thị và dữ kiện”. Phát biểu nào sau đây đúng?
\choice
{Mọi quá trình nhiệt học đều không phụ thuộc điều kiện thực hiện.}
{\True Khí nhận nhiệt và thực hiện công có ΔU=Q-A_do_khí_thực_hiện.}
{Khí nhận nhiệt thì nội năng luôn tăng đúng bằng nhiệt lượng nhận được.}
{Nhiệt lượng và nhiệt độ là hai đại lượng hoàn toàn giống nhau.}
\loigiai{
Phát biểu đúng là: Khí nhận nhiệt và thực hiện công có ΔU=Q-A_do_khí_thực_hiện. Các phương án còn lại trái với mô hình, định nghĩa hoặc hệ thức nhiệt học tương ứng.
}
\end{ex}

% ID: L12C1B2-28-TN
% Mức: VD
% ID: L12C1B2-28-TN
% ID: L12C1B2-28-TN
% Mức: VD
% ID: L12C1B2-28-TN
% ID: L12C1B2-28-TN
% Mức: VD
% ID: L12C2B13-76-TN
%=== Câu 73 ===%
\begin{ex}
% ID: L12C2B13-76-TN
% Mức: VD
Câu 7 thuộc dạng “Nhận diện quá trình nhiệt động từ đồ thị và dữ kiện”. Phát biểu nào sau đây đúng?
\choice
{Mọi quá trình nhiệt học đều không phụ thuộc điều kiện thực hiện.}
{\True Khí nhận nhiệt và thực hiện công có ΔU=Q-A_do_khí_thực_hiện.}
{Khí nhận nhiệt thì nội năng luôn tăng đúng bằng nhiệt lượng nhận được.}
{Nhiệt lượng và nhiệt độ là hai đại lượng hoàn toàn giống nhau.}
\loigiai{
Phát biểu đúng là: Khí nhận nhiệt và thực hiện công có ΔU=Q-A_do_khí_thực_hiện. Các phương án còn lại trái với mô hình, định nghĩa hoặc hệ thức nhiệt học tương ứng.
}
\end{ex}

% ID: L12C1B2-29-TN
% Mức: VDC
% ID: L12C1B2-29-TN
% ID: L12C1B2-29-TN
% Mức: VDC
% ID: L12C1B2-29-TN
% ID: L12C1B2-29-TN
% Mức: VDC
% ID: L12C2B13-77-TN
%=== Câu 75 ===%
\begin{ex}
% ID: L12C2B13-77-TN
% Mức: VDC
Câu 4 thuộc dạng “Nhận diện quá trình nhiệt động từ đồ thị và dữ kiện”. Phát biểu nào sau đây đúng?
\choice
{\True Khí nhận nhiệt và thực hiện công có ΔU=Q-A_do_khí_thực_hiện.}
{Khí nhận nhiệt thì nội năng luôn tăng đúng bằng nhiệt lượng nhận được.}
{Nhiệt lượng và nhiệt độ là hai đại lượng hoàn toàn giống nhau.}
{Mọi quá trình nhiệt học đều không phụ thuộc điều kiện thực hiện.}
\loigiai{
Phát biểu đúng là: Khí nhận nhiệt và thực hiện công có ΔU=Q-A_do_khí_thực_hiện. Các phương án còn lại trái với mô hình, định nghĩa hoặc hệ thức nhiệt học tương ứng.
}
\end{ex}

% ID: L12C1B2-30-TN
% Mức: VDC
% ID: L12C1B2-30-TN
% ID: L12C1B2-30-TN
% Mức: VDC
% ID: L12C1B2-30-TN
% ID: L12C1B2-30-TN
% Mức: VDC
% ID: L12C2B13-78-TN
%=== Câu 77 ===%
\begin{ex}
% ID: L12C2B13-78-TN
% Mức: VDC
Câu 8 thuộc dạng “Nhận diện quá trình nhiệt động từ đồ thị và dữ kiện”. Phát biểu nào sau đây đúng?
\choice
{\True Khí nhận nhiệt và thực hiện công có ΔU=Q-A_do_khí_thực_hiện.}
{Khí nhận nhiệt thì nội năng luôn tăng đúng bằng nhiệt lượng nhận được.}
{Nhiệt lượng và nhiệt độ là hai đại lượng hoàn toàn giống nhau.}
{Mọi quá trình nhiệt học đều không phụ thuộc điều kiện thực hiện.}
\loigiai{
Phát biểu đúng là: Khí nhận nhiệt và thực hiện công có ΔU=Q-A_do_khí_thực_hiện. Các phương án còn lại trái với mô hình, định nghĩa hoặc hệ thức nhiệt học tương ứng.
}
\end{ex}
